\clearpage
\setcounter{section}{6}
\setcounter{subsection}{4}
\setcounter{equation}{17}

Nach dem Debye Modell für die spezifische Wärme gilt für die Energie
\begin{equation}
	\label{eq:II.6.18}
	E_{n} = \hbar  \omega \left( n + \frac{1}{2} \right) = n \hbar  \omega.
\end{equation}	

\paragraph{Anmerkung} Hierdurch wird die Nullpunktsenergie vernachlässigt. Sie spielt fúr die Berechnung der Wärmekapazität keine Rolle da
$$
C_{V} = \frac{\mathrm{d}U}{\mathrm{d}T} 
$$ 
gilt.\\
Es ergibt sich mit \ref{eq:II.6.11}
\begin{equation}
	\label{eq:II.6.19}
	U = \int \mathrm{d} \omega D(\omega) \left< n(\omega) \right> \hbar  \omega = \int_{0}^{\omega_{D}} \mathrm{d} \omega \underbrace{ \frac{3 V \omega^2}{2 \pi^2 v_{\mathrm{g}}^3}}_{D(\omega)} \underbrace{ \frac{\hbar  \omega}{e^{\hbar  \omega /k_{\mathrm{B}} T} - 1}}_{\left<n(\omega) \right> \hbar  \omega}
\end{equation}
und mit
$$
x = \frac{\hbar  \omega}{k_{\mathrm{B}}T} \quad \text{und}\quad x_{D} = \frac{\hbar \omega_{D}}{k_{\mathrm{B}}T} = \frac{\Theta}{T}	
$$ 
folgt
\begin{equation}
	\label{eq:II.6.20}
	U = 9 N k_{\mathrm{B}}T \left( \frac{T}{\Theta} \right) ^3 
	\int_{0}^{x_{D}} \mathrm{d}x \frac{x^3}{e^{x} -1}
\end{equation}
und
\begin{equation}
	\label{eq:II.6.21}
	C_{V} = \left( \frac{\mathrm{d}U}{\mathrm{d}T}\right)_{V} = 9Nk_{B} \left( \frac{T}{\Theta} \right) ^3 \int_{0}^{x_{D}} \mathrm{d} x \frac{x^4 e^{x}}{\left( e^{x} -1  \right) ^2}
\end{equation}	

\begin{figure}[H]
    \centering
    \incfig{vl24_abbildung_1}
    \label{fig:vl24_abbildung_1}
\end{figure}

\paragraph{Diskussion}\,\\
\paragraph{Klassisches Resultat}\,\\ 
Für $T\gg\Theta$, $x_{D} \ll 1$, $e^{x} \approx 1$, $e^{x}-1=x$ gilt
		$$
		\int_{0}^{x_{D} x^2 \mathrm{d}x} = \frac{x_{D}^3}{3} = \frac{1}{3} \left( \frac{\Theta}{T} \right) ^3
		$$ 
		und es folgt das klassische Resultat
		\begin{equation}
			\label{eq:II.6.22}
			C_{V} = 3R
		\end{equation}


\paragraph{Debye'sche $T^3$ Gesetz}\,\\ 
Für $T \ll \Theta$, $x_{D} \gg 1$, Integrationsgrenze $\to $ $\infty$ 
		$$
		\int_{0}^{\infty} \frac{x^{4} e^{x}}{\left( e^{x} -1 \right) ^2} \mathrm{d}x = \frac{4\pi^{4}}{15}
		$$ 
		Das Debye'sche $T^3$ Gesetz
		\begin{equation}
			\label{eq:II.6.23}
			C_{V}= \frac{12 \pi^{4}R}{5 \Theta^3} T^3 \quad \left( T\le \num{0,1} \Theta \right) 
		\end{equation}


\paragraph{Experimentelle Bestimmung}\,\\ 
Bestimmung von $C_{V}(T)$ $\to $ Anpassung mit Theorie $\implies$ $\Theta$
\begin{table}[htpb]
	\centering
	\caption{caption}
	\label{tab:label}
	\begin{tabular}{|c|c|c|}
	\hline
	 Stoff & $\Theta$ & $\omega_{D}$ \\
	 \hline
	 \ce{Na} & \SI{158}{\kelvin} & \SI{2 e13}{\hertz} \\
	 \hling
	 \ce{Ca} & \SI{343}{\kelvin} & \SI{4,5 e 13}{\hertz}\\
	\hline
	\end{tabular}
\end{table}

\paragraph{Anpassung}\,\\ 
Bessere Anpassung and die gesamte Kurve $C_{V}(T)$ mit Dispersionsbeziehung $\omega(k)$.

\paragraph{Abweichung}\,\\
Abweichung zum $T^3$-Gesetz: Beitrag der Elektronen zur Wärmekapazität
\paragraph{Proportionalität}\,\\ 
Es gilt
$$
N_{\text{Ph}} \sim \begin{cases}
	T^3 &\text{für} T\ll \Theta_{D} \\
	T & \text{für} T\gg \Theta_{D}
\end{cases}
$$ 
was für bei verschiedenen Phänomenen eine Rolle spielt.

\paragraph{$T=\Theta$}\,\\ 
Für $T=\Theta$ ist $U_{\text{th}}\simeq 2 U_{0}$, wobei $U_0$ der thermischen Energie der Nullpunktsschwankung entspricht.

\paragraph{Klassische und QM Betrachtung}\,\\ 
Bei $T<\Theta_{D}$ frieren Moden aus und das resultierende Verhalten ist nur quantenmechanisch erklärbar. Bei $T > \Theta_{D}$ werden alle Moden angeregt und das Verhalten ist klassich erklärbar \textcolor{red}{war nicht in der Vorlesung, Satz ist wsl falsch}.

\subsection{Wärmeleitung des Kristallgitters}
Für die Stromdichte $j$, also die Energie die pro Zeiteinheit durch ein Flächeneinheit hindurch tritt, gilt
\begin{equation}
	\label{eq:II.6.24}
	\frac{\dot{Q}}{F} = j = - k \frac{\mathrm{d}T}{\mathrm{d}x}
\end{equation}
mit der Fläche $F$, der Energie pro Zeiteinheit $\dot{Q}$ und dem Wärmeleitvermögen $k$.

\paragraph{Messung}\,\\ 
Die Messung erfolgt mit der in Abbildung \ref{fig:vl24_abbildung_2} schematisch abgebildeten Apparatur.

\begin{figure}[H]
    \centering
    \def\svgwidth{0.59\columnwidth}
    \subimport{figures/}{vl24_abbildung_2.pdf_tex}
    \label{fig:vl24_abbildung_2}
\end{figure}

Dabei wird auf einen stationären Zustand gewartet, bei welchem
$$
\dot{Q} = U I = - k F \frac{\Delta T}{l}
$$ 
gilt.

\paragraph{Wärmeleitung (Phononen als idealer Glas)}\,\\ 
Die statistische Natur des Leitungsprozesses bringt den Temperaturgradienten in die Gleichung \ref{eq:II.6.24}, Energie diffundiert durch die Probe. \\
Phänomenologisch :
\begin{itemize}
	\item mittlere Lebensdauer $\tau$
	\item freie Weglänge $\lambda = v_{\text{g}} \tau$ (Umklappprozesse)
	\item Energie $\hbar  \omega$
\end{itemize}
\begin{figure}[H]
    \centering
    \def\svgwidth{0.369\columnwidth}
    \subimport{figures/}{vl24_abbildung_3.pdf_tex}
    \label{fig:vl24_abbildung_3}
\end{figure}

In Stromrichte in Richtung einer Koordinatenachse berechnet sich mit
$$
j_{x} = \underbrace{\underbrace{- \frac{1}{V} \Delta \left< n \right>}_{\text{Teilchendichte}} \hbar \omega}_{\text{Energiedichte}} v_{\text{g}x} 
$$  
mit
$$
\Delta \left<n \right> = \frac{\mathrm{d}\left<n \right>}{\mathrm{d}x} \Delta x \quad \text{und}\quad \Delta x = \lambda = v_{\text{g}} \tau 
$$
mit der freien Phononenweglänge $\lambda$, welche sich auf Umklappprozesse bezieht.
$$
\begin{aligned}
	j_{x} &= - \frac{1}{V} \hbar  \omega \frac{\mathrm{d}\left< n \right>}{\mathrm{d}x} v_{\text{g}x}^2 \tau  \\
	&= \frac{1}{V} \hbar \omega \frac{\mathrm{d} \left< n \right>}{\mathrm{d}T} \frac{\mathrm{d}T}{\mathrm{d}x} v_{\text{g}x}^2 \tau   
\end{aligned}
$$ 
Genauer: Summation über alle Schwingungsmoden: $k_{j}, \omega\left( k_{j} \right) $ der $j$ Äste
\begin{equation}
	\label{eq:II.}
	j_{x} = - \underbrace{\frac{1}{V} \sum_{\Vec{k}, j}^{} \hbar \omega_{j} \left( k_{j} \right) \tau_{j} (k_{j}) v_{\text{g},x,j}^2 (k_{j}) \frac{\mathrm{d} \left<n_{j} (k_{j}) \right>}{\mathrm{d}T}}_{k} \frac{\mathrm{d}T}{\mathrm{d}x} 
\end{equation}
Vergleich mit \ref{eq:II.6.24}, für isotrope Materialien
$$
\left<v_{x}^2 \right> = \frac{1}{3} v_{\text{g}}^2
$$ 
und mit $\lambda = v_{\text{g}} \tau $ folgt
\begin{equation}
	\label{eq:II.6.26}
	k = \frac{1}{3} \sum_{\Vec{k}, j}^{}  v_{\text{g},j} \lambda_{j} \underbrace{\frac{1}{V} \frac{\partial }{\partial T} \varepsilon(\omega_{j}, T)}_{C_{V,j}}
\end{equation}
Für die weitere Diskussion berücksichtigt man einen Phononenast (Debye-Näherung) und wendet die \textbf{dominante Phononennäherung} an, d.h. das Frequenzspektrum wird vernachlässigt und man betrachtet nur Phononen, die am meisten zum Wärmetransport beitragen ($\hbar \overline{\omega} \simeq p k_{\mathrm{B}} T$ wobei $1 \le p \le 3$)

\begin{equation}
	\label{eq:II.6.27}
	k = \frac{1}{3} v_{\mathrm{g}} \lambda C_{V}
\end{equation}

\paragraph{Phononen machen Stöße}\,\\ 
\begin{enumerate}
	\item mit den Wänden 
	\item mit Gitterfehlern
	\item mit anderen Phononen
\end{enumerate}
$$
\frac{1}{\lambda} = \frac{1}{\lambda_{\text{a)}}} + \frac{1}{\lambda_{\text{b)}}} + \frac{1}{\lambda_{\text{c)}}}
$$ 
Prozess, der zur stärksten Streuung führt bestimmt, $\lambda$.\\
bei harmonischen Potential besteht keine Phonon-Phonon-Wechselwirkung und somit keine Wärmeleitung, wohingegen bei anhorminischen Potential Phonon-Phonon-Wechselwirkung besteht.

\begin{figure}[H]
    \centering
    \incfig{vl24_abbildung_4}
    \label{fig:vl24_abbildung_4}
\end{figure}

\paragraph{Normalprozess: N-Prozess}\,\\ 
Hier hat man einen energiefluss von links nach rechts, kein Temperaturgradient
$$
\begin{aligned}
	\hbar \Vec{k}_{1} + \hbar \Vec{k}_{2} &= \hbar \Vec{k}_{3} \\
	\hbar \omega_{1} + \hbar \omega_{2} &= \hbar \omega_{3} 
\end{aligned}
$$ 

\paragraph{Umklappprozesse: U-Prozess}\,\\ 
Ein Temperaturgradient entsteht!
\begin{figure}[H]
    \centering
    \incfig{vl24_abbildung_5}
    \label{fig:vl24_abbildung_5}
\end{figure}

$$
\hbar \Vec{k}_{1} + \hbar \Vec{k}_{2} &=  \hbar \Vec{k}_{3} + \hbar \Vec{G} \\
\hbar  \omega_1 + \hbar \omega_{2} &= \hbar \omega_3
$$ 
U-Prozesse können nur die phononen machen für die gilt
\begin{enumerate}
	\item 
		$$
		k_{1,2} \ge \frac{\pi/a}{2} = \frac{k_{\text{max}}}{2}
		$$ 
	\item $\omega = v_{\text{g}} k$ (Debye Modell) $\omega_{\text{max}}= \omega_{D}$ 
		\begin{equation}
			\label{eq:II.6.28}
			\hbar \omega_{1 /2} = \frac{\hbar  \omega_{\text{max}}}{2} = \frac{k_{\mathrm{B}} \Theta}{2}
		\end{equation}
\end{enumerate}

\paragraph{Abschätzung für $k$ für einen Isolator}\,\\ 
Die Phononengeschwindigkeit sei konstant mit $v_{\text{g}} = \SI{5 e5}{\centi \meter \per \second}$. Es gilt $C_{V} = \SI{24,9}{\joule \mol^{-1} \kelvin^{-1}}$ $T \gg \Theta$, und $C_{V} \sim T^3$ für $T\ll\Theta$. Die Wechselwirkung mit Phononen ist proportional zur mittleren Phononenzahl
\begin{equation}
	\label{eq:II.6.29}
	\lambda = \lambda(T) \sim \frac{1}{\left< w \right>} \sim e^{\Theta / 2T} -1
\end{equation}

$$
\left< n \right> = \frac{1}{e^{\hbar \omega / k_{\mathrm{B}}T}} \simeq \frac{1}{e^{\hbar  \omega_{\text{max} / 2 k_{\mathrm{B}}T}}-1} = \frac{1}{e^{k_{\mathrm{B}} \Theta /2 k_{B}T} - 1}
$$ 
Für
\begin{enumerate}
	\item $T\gg\Theta$ gilt
		$$
		\begin{aligned}
			C_{V} &= \text{konst.}(T) \\
			\lambda &\sim \frac{\Theta}{2T} \sim \frac{1}{T} \\
			k &\sim \frac{1}{T} \\
		\end{aligned}
		$$ 
	\item $T \le \Theta$ gilt
		$$
		\begin{aligned}
			C_{V} &\simeq \text{konst}(T) \\
			\lambda &\sim e^{\Theta /2T}\\
			k &\sim e^{\Theta/2T}\\
		\end{aligned}
		$$ 
	\item $T\ll\Theta$ gilt
		$$
		\begin{aligned}
			C_{V} &\sim T^3 \\
			\lambda &\sim e^{\Theta /2T} > d = \, \text{Probedimension}\\
			k &\sim \mathrm{d} T^3 \, \text{(Casimir-Bereich)}
		\end{aligned}
		$$ 
\end{enumerate}

\begin{figure}[H]
    \centering
    \incfig{vl24_abbildung_6}
    \label{fig:vl24_abbildung_6}
\end{figure}

Unreine Kristalle: $a$ mittlerer Abstand zwischen 2 Fehlstellen, $T\ll \Theta$ $\lambda \sim a$ $k \sim aT^3$, $a<d$
